\documentclass[11pt]{article}
\usepackage{amsmath,amssymb}

\title{Progressive Growing GANs for High-Fidelity Image Synthesis}
\author{D. Karlsson \and A. Mehta \and P. Lefevre}
\date{}

\begin{document}
\maketitle

\begin{abstract}
We describe a training procedure that grows generator and discriminator
symmetrically from $4\times4$ to $1024\times1024$ resolution by smoothly
fading in new layers. The schedule decouples the difficulty of capturing
coarse structure from the difficulty of resolving fine texture, and yields
the first generative model to produce uncurated megapixel face images that
are routinely confused with photographs in side-by-side trials.
\end{abstract}

\section{Motivation}
End-to-end GAN training at high resolution is unstable: the discriminator
finds high-frequency artifacts long before the generator has settled the
underlying composition, and the resulting gradient is noisy and adversarial
in the worst sense. Progressive growing sidesteps this by giving the model
time to converge at each scale before increasing it.

\section{Method}
We parameterise the generator $G$ as a sequence of upsampling blocks and the
discriminator $D$ as the symmetric stack. New layers are introduced with a
fade weight $\alpha \in [0, 1]$:
\begin{equation}
G_\text{out} = \alpha \cdot G_\text{new}(z) + (1 - \alpha) \cdot
\mathrm{upsample}(G_\text{old}(z)).
\end{equation}
$\alpha$ ramps linearly from $0$ to $1$ over a fixed number of iterations,
after which the lower-resolution path is discarded. Two further stabilisers
matter: equalised learning rates (rescaling weights at runtime so all layers
have effectively unit variance updates) and pixelwise feature normalisation
in the generator.

\section{Loss}
We use the Wasserstein-GP loss for its calmer gradient landscape:
\begin{equation}
\mathcal{L} = \mathbb{E}_{\tilde{x} \sim p_g}[D(\tilde{x})] -
              \mathbb{E}_{x \sim p_r}[D(x)] +
              \lambda \, \mathbb{E}_{\hat{x}}\!\left[(\|\nabla_{\hat{x}} D(\hat{x})\|_2 - 1)^2 \right],
\end{equation}
with $\lambda = 10$ and $\hat{x}$ sampled along straight lines between real
and fake batches.

\section{Implementation}
Eight V100 GPUs train each resolution stage for 600k images. Batch size halves
with each doubling: 256 at $4\times4$, down to 8 at $1024\times1024$. The full
schedule on FFHQ takes about nine days of wall-clock time.

\section{Evaluation}
\begin{tabular}{lcc}
\hline
Resolution & FID & Hours/stage \\
\hline
$128^2$  & 4.1  & 8 \\
$256^2$  & 5.6  & 14 \\
$512^2$  & 6.9  & 32 \\
$1024^2$ & 7.7  & 96 \\
\hline
\end{tabular}

FID degrades only mildly with resolution, contrasting with the typical
collapse seen when models are trained at the target resolution from scratch.
A user study with $N=200$ participants finds 47\% accuracy at distinguishing
$1024^2$ samples from real FFHQ photographs --- statistically indistinguishable
from chance.

\section{Failure Modes}
Mode collapse on rare attributes (eyeglasses, hats) persists; truncation tricks
trade fidelity for diversity at sampling time. We also see characteristic
checkerboard patterns when the upsampling kernel size and stride disagree ---
a known consequence of transpose convolution that nearest-neighbour
upsampling avoids.

\section{Discussion}
Progressive growing was eventually subsumed by style-based generators, but its
core insight --- that curriculum on resolution is curriculum on frequency
content --- carries over directly. The right way to think about generative
training is in the spectrum, not in the pixel grid.

\bibliographystyle{plain}
\begin{thebibliography}{9}
\bibitem{karras2018progressive} T. Karras et al. \emph{Progressive Growing of GANs.} ICLR, 2018.
\bibitem{gulrajani2017improved} I. Gulrajani et al. \emph{Improved Training of Wasserstein GANs.} NeurIPS, 2017.
\end{thebibliography}

\end{document}
